% !TEX root = IE3_expH_report.tex
\documentclass[lualatex,a4paper,10pt]{jlreq}
% Shared LuaLaTeX preamble.
\usepackage{luatexja-fontspec}
\setmainfont{TeX Gyre Termes}
\setsansfont{TeX Gyre Heros}
\setmonofont{Latin Modern Mono}
\setmainjfont[BoldFont=HaranoAjiGothic-Medium]{HaranoAjiMincho-Regular}
\setsansjfont[BoldFont=HaranoAjiGothic-Bold]{HaranoAjiGothic-Medium}
\usepackage[top=22mm,bottom=22mm,left=23mm,right=23mm]{geometry}
\usepackage{amsmath,amssymb,booktabs,array,longtable,tabularx}
\usepackage{graphicx,xcolor,tikz,circuitikz}
\usetikzlibrary{arrows.meta,positioning,calc}
\usepackage{enumitem,fancyhdr,listings,needspace}
\usepackage[unicode,colorlinks=true,linkcolor=labblue,urlcolor=labteal]{hyperref}
\usepackage{bookmark}
\definecolor{labblue}{RGB}{30,70,112}
\definecolor{labteal}{RGB}{0,100,105}
\definecolor{labgray}{gray}{0.95}
\ModifyHeading{section}{font={\large\sffamily\bfseries\color{labblue}}}
\ModifyHeading{subsection}{font={\normalsize\sffamily\bfseries\color{labteal}}}
\setlength{\parindent}{1\zw}
\setlength{\parskip}{3pt}
\setlength{\emergencystretch}{2em}
\renewcommand{\arraystretch}{1.35}
\setlist{itemsep=3pt,topsep=4pt,leftmargin=2\zw}
\newcolumntype{Y}{>{\raggedright\arraybackslash}X}
\newcolumntype{C}[1]{>{\centering\arraybackslash}p{#1}}
\pagestyle{fancy}
\fancyhf{}
\fancyhead[L]{\small\color{labblue} IE3 後期・電子工学実験}
\fancyhead[R]{\small テーマ\LabCode}
\fancyfoot[C]{\thepage}
\setlength{\headheight}{16pt}
\DeclareOldFontCommand{\rm}{\normalfont\rmfamily}{\mathrm}
\lstset{basicstyle=\small\ttfamily,columns=fullflexible,keepspaces=true,
 breaklines=true,frame=single,backgroundcolor=\color{labgray},
 numbers=left,numberstyle=\tiny,showstringspaces=false,aboveskip=8pt,belowskip=8pt}
\newcommand{\blank}{\rule{22mm}{0.3pt}}
\newcommand{\answerline}{\par\noindent\rule{\linewidth}{0.25pt}\par}
\newcommand{\writing}[1]{\par\Needspace{\dimexpr#1+5mm\relax}\noindent%
 \fcolorbox{labblue!35}{white}{\parbox[c][#1][t]{\dimexpr\linewidth-2\fboxsep-2\fboxrule\relax}{\vspace{2mm}\noindent\strut\hfill\par}}\par}
\newcommand{\hint}[1]{\par\smallskip\noindent{\sffamily\bfseries\color{labteal}記述の視点：}#1\par}
\newcommand{\note}[1]{\par\smallskip\noindent\fcolorbox{labblue!45}{labblue!5}{%
 \parbox{\dimexpr\linewidth-2\fboxsep-2\fboxrule\relax}{#1}}\par\smallskip}
\newcommand{\eqerror}{\varepsilon=\frac{x_{\mathrm{meas}}-x_{\mathrm{th}}}{x_{\mathrm{th}}}\times100\ \%}
\newcommand{\LabSource}{徳山工業高等専門学校 情報電子工学科，
 『2026年度 電子工学実験（3年後期）テキスト』，テーマ\LabCode，
 \expandafter\path\expandafter{\SourcePdf}。}
\newcommand{\reporttitle}{
 \noindent{\small\sffamily\color{labteal}実験報告書・記入ガイド}\par
 \noindent{\LARGE\sffamily\bfseries \LabCode\quad\LabTitle}\par\smallskip
 \noindent 氏名：\blank\quad 班：\blank\quad 実験日：\blank\par
 \note{目的・原理・手順を確認し、実測値と自分の考察を記入する。
 考察のガイドは、検討する事象・使う測定結果・記述の方針を示す。}
}
\newcommand{\prelabtitle}{
 \noindent{\small\sffamily\color{labteal}事前学習・前提知識プリント}\par
 \noindent{\LARGE\sffamily\bfseries \LabCode\quad\LabTitle}\par\smallskip
 \noindent 氏名：\blank\quad 班：\blank\quad 予習日：\blank\par
 \note{用語、回路の動き、計算、測定表への記入の順に読む。
 例題の値は説明用であり、実験結果ではない。}
}
\newcommand{\tocrow}[3]{\hyperref[#1]{#2}& #3 &\pageref*{#1}\\[2mm]}
\newcommand{\documentmap}[1]{
 \section*{目次と概要}
 \noindent 青色の項目をクリックすると該当箇所へ移動する。\par
 {\small\begin{longtable}{@{}p{0.29\linewidth}p{0.61\linewidth}r@{}}
 \toprule {\color{labblue}\bfseries 項目} & {\color{labblue}\bfseries 読むと分かること} & 頁\\\midrule
 \endhead #1\bottomrule
 \end{longtable}}
 \clearpage
}
\newenvironment{terms}{\par\smallskip\begin{longtable}{@{}p{0.23\linewidth}p{0.73\linewidth}@{}}
 \toprule {\color{labteal}\bfseries 用語}&{\color{labteal}\bfseries 意味・回路での役割}\\\midrule\endhead}
 {\bottomrule\end{longtable}}
\newcommand{\term}[2]{\textbf{#1}&#2\\[2mm]}
\newcommand{\guidepoint}[4]{
 \Needspace{32mm}\subsection{#1}
 \noindent{\bfseries\color{labteal}考える事象：}#2\par
 \noindent{\bfseries\color{labteal}参照する結果：}#3\par
 \noindent{\bfseries\color{labteal}記述の方針：}#4\par
}
\newcommand{\reportclosing}[1]{
 \clearpage\section{結論のまとめ方}\label{sec:closing}
 \hint{#1}
 \ConclusionGuide
 \begin{enumerate}
 \item 目的に対応する最も重要な実測結果を、測定条件と数値を添えて述べる。
 \item 原理と一致した点・一致しなかった点を示す。原因は確認済みの事実と推測を分ける。
 \item 実施範囲で言えることと、未確認の条件を一文ずつまとめる。
 \end{enumerate}
 \note{書き出しの型：「○○の条件で△△を測定したところ、□□となった。
 これは◇◇と整合する。ただし××は確認しておらず、一般化には追加測定が必要である。」
 空欄は自分の実測値と根拠で埋める。}
 \noindent 結論（実験後に記入）：\writing{30mm}
 \noindent 改善案・次に確認したい条件：\writing{16mm}
 \section*{参考文献}\label{sec:references}
 \begin{enumerate}
 \item \LabSource
 \item 使用したデータシート・書籍・ウェブ資料（著者、題名、該当箇所、URL、参照日）を追加する。
 \end{enumerate}
}
\newcommand{\prelabclosing}{
 \section*{準備・出典}\label{sec:prelabsource}
 \begin{itemize}[nosep]
 \item 資料、筆記具、記録用紙、電卓と、実験条件・機器設定の記録欄。
 \item 確認問題の解答と、分からない用語・回路点への具体的な質問。
 \end{itemize}
 {\small\noindent\LabSource\par}
}
\newcommand{\simpleflow}[1]{
 \begin{center}
 \begin{tikzpicture}[node distance=5mm and 5mm,
 every node/.style={font=\small},box/.style={draw=labblue,fill=labblue!4,rounded corners=1mm,
 align=center,minimum height=10mm,text width=30mm}]
 #1
 \end{tikzpicture}
 \end{center}
}

\newcommand{\LabCode}{H}
\newcommand{\LabTitle}{PICを用いたフルカラーLED制御}
\newcommand{\SourcePdf}{IE3_expH_A4.pdf}
\newcommand{\ConclusionGuide}{負論理の色対応とLOW割合の実測を中心に、点灯・調光・混色を実現できた範囲をまとめる。
電流を実測していない場合は抵抗による推定と明記し、明るさの評価方法も添える。\par}
\begin{document}
\reporttitle
\documentmap{
\tocrow{sec:auto1}{1 目的}{負論理GPIO、電流制限、RGBのPWM調光を実験で確かめる。}
\tocrow{sec:auto2}{2 原理}{負論理GPIO、電流制限、RGBのPWM調光の仕組みと成立条件を整理する。}
\tocrow{sec:auto3}{3 装置・設定}{機器・定数・測定点を確認し、資料の順序で操作する。}
\tocrow{sec:auto4}{4 方法}{機器・定数・測定点を確認し、資料の順序で操作する。}
\tocrow{sec:auto5}{5 結果}{ポートと色、LOW時間、混色を表や図に記録する。}
\tocrow{sec:auto6}{6 プログラム添付}{この項目の仕組みと実験で確認すべき点を整理する。}
\tocrow{sec:auto7}{7 考察：結果を根拠に説明する}{ポートと色、LOW時間、混色を根拠に、比較と原因の検討を進める。}
\tocrow{sec:closing}{8 結論のまとめ方}{主要な結果、成立範囲、未確認の条件を短くまとめる。}
\tocrow{sec:references}{参考文献}{使用した資料と外部の根拠を示す。}
}

\section{目的}\label{sec:auto1}
PIC12F683でフルカラーLEDの点灯・混色・PWM調光を行い、GPIOの論理、電流制限、プログラムと出力波形の関係を理解する。
\section{原理}\label{sec:auto2}
共通アノード型RGB LEDは、各カソードにつながる出力をLOWにすると点灯する。
$GP0$が赤、$GP1$が緑、$GP2$が青である。
抵抗の近似値は$R=(V_{CC}-V_F)/I_F$。
PWMでは点灯時間の割合$D_{\rm on}=t_{\rm LOW}/T$を変えて平均電流を制御する。高レベルのデューティ比とは逆になる。
\section{装置・設定}\label{sec:auto3}
PIC12F683、LED実験ボード、PC、資料指定のMPLAB/コンパイラ、PICkit3、電源、オシロスコープを使用する。
資料のプロジェクト設定を読み込み、書込み時の電源設定と動作時5 Vを区別する。
内蔵発振器4 MHz、GPIO出力方向、ANSEL、比較器を設定する。GP3は入力専用である。
実際のLED型番、抵抗値、開発環境、電源設定：
\writing{16mm}
\section{方法}\label{sec:auto4}
\begin{enumerate}
\item 電源OFFで装置を接続し、サンプル1を書き込む。赤・緑・青が2秒ごとに順次点灯することを確認する。
\item GP2,GP1,GP0の全8組を出力し、資料の表H-2を観測結果で埋める。
\item サンプル2の赤の5段階調光を確認する。青を1秒ごとに明るくし、最大後は1秒ごとに暗くするプログラムを作成する。
\item 最大輝度と3段階目でGP2--GND間の電圧を観測し、周期とLOW時間を測定する。
\item 赤$\to$黄$\to$緑$\to$水色$\to$青$\to$紫$\to$赤と変化するプログラムを作成する。各RGB色の間を約5段階で変化させる。
\end{enumerate}
LEDを直接見続けない。電流制限抵抗を省略せず、資料の機器操作に従う。
\clearpage
\section{結果}\label{sec:auto5}
表1：ポート状態と混色。理論予想と観測を分ける。
\par\noindent\begin{tabularx}{\linewidth}{YYYYY}
\toprule GP2 & GP1 & GP0 & 予想色 & 観測色\\\midrule
0&0&0&&\\0&0&1&&\\0&1&0&&\\0&1&1&&\\
1&0&0&&\\1&0&1&&\\1&1&0&&\\1&1&1&&\\\bottomrule
\end{tabularx}
表2：青のPWM調光。
\par\noindent\begin{tabularx}{\linewidth}{YYYY}
\toprule 段階 & 周期/ms & LOW時間/ms & 点灯割合/\%\\\midrule
消灯&&&\\1&&&\\2&&&\\3&&&\\4&&&\\最大&&&\\\bottomrule
\end{tabularx}
図1：最大と3段階目のGP2波形。電圧・時間・点灯区間を示す。
\writing{38mm}
図2：混色プログラムの流れ、RGB点灯割合の変化。
\writing{30mm}
\clearpage
\section{プログラム添付}\label{sec:auto6}
基本課題2および応用課題のソースを添付し、初期設定、PWM一周期、段階の保持時間、色の遷移を説明する。
\writing{22mm}
\clearpage\section{考察：結果を根拠に説明する}\label{sec:auto7}
\guidepoint{ポート状態と混色}
{全8状態の発光色が負論理の予想に対応するか。}
{表1のGP2 GP1 GP0、予想色と観測色、RGBの配線。}
{各状態の点灯色をビットから説明する。見え方が違う場合は色ごとの電流や光量を候補とし、電圧測定だけで光量を確定しない。}
\guidepoint{PWMの波形と調光}
{最大と3段階目の出力が、設計したLOW割合になるか。}
{表2、GP2--GND波形、周期・LOW時間、段階保持の観察。}
{HIGHデューティを点灯割合に取り違えない。最大時に連続LOWならその状態を示し、一秒の段階周期とPWM周期を別々に説明する。}
\guidepoint{電流制限と端子の動作条件}
{抵抗値がLED・GPIOの条件に対して適切か。}
{実装抵抗、電源、LED型番・順方向電圧、公式データシート。}
{指定条件の理論電流と実装の推定電流を分ける。定格・出力電圧の制約を確認し、計算だけなら「実測電流」と呼ばない。}
\guidepoint{色の連続変化とプログラム}
{赤・緑・青の間で約5段階の変化が期待通りに見えるか。}
{ソース、色ごとの割合の時系列、実際の順序と段階の観察。}
{色順と割合の更新規則を示し、段階数・間隔・見え方の対応を説明する。人の感覚による評価と電気波形による確認を区別する。}
\subsection{資料の課題・追加検討}
\begin{enumerate}
\item LEDの用途と型番・価格を調べ、販売元、調査日、個数条件、URLを示す。価格は調査した値を記す。
\item 赤の$V_F=2.0\,\mathrm{V}$、緑・青$3.4\,\mathrm{V}$、電源5 V、20 mAという設計条件で抵抗を計算する。実装抵抗から実際の想定電流も求める。
\item GPIOのソース／シンク動作を説明し、データシートの出力電流制限と実装条件を確認する。
\item 定電流ダイオードの特徴、必要電圧、抵抗による電流制限との違いを説明する。
\item PWMの周波数・点灯割合と観測した明るさの関係を説明する。色ごとの明るさの差も検討する。
\end{enumerate}
\reportclosing{負論理の点灯とPWM波形、混色プログラムの動作をまとめる。}
\end{document}
